\documentclass[12pt,a4paper]{article}
\usepackage[utf8]{inputenc}
\usepackage{amsmath,amssymb,amsthm}
\usepackage{mathtools}
\usepackage{geometry}
\geometry{margin=1in}

\newtheorem{theorem}{Theorem}[section]
\newtheorem{definition}[theorem]{Definition}
\newtheorem{remark}[theorem]{Remark}
\newtheorem{conjecture}{Conjecture}[section]

\title{Original Contributions from a Greedy Harmonic Approach to the Riemann Zeta Function}
\author{}
\date{}

\begin{document}

\maketitle

\begin{abstract}
This article summarises several original insights that arose from an investigation into the connection between greedy harmonic reconstruction and the Dirichlet eta function. The contributions include: (1) a rigorous restriction of a functional equation for modular tensor categories to the case of lattice MTCs; (2) a greedy algorithm for reconstructing sine from harmonic fractions; (3) a reformulation of the Riemann Hypothesis (RH) in terms of a greedy condition on the partial sums of the eta function; (4) a heuristic argument identifying the critical exponent $1/2$ via backward propagation from the first term; (5) a research roadmap using modern analytic number theory; (6) a clarification of why a simple density argument fails; and (7) the concept of inverse parametric eta. No proof of RH is claimed, but these ideas provide a new perspective and a testable reformulation.
\end{abstract}

\section{Introduction}

The Riemann zeta function $\zeta(s)$ and its Dirichlet eta variant $\eta(s)=\sum_{n=1}^\infty (-1)^{n-1}n^{-s}$ have been studied for over a century. The Riemann Hypothesis (RH) states that all non‑trivial zeros of $\zeta(s)$ lie on the line $\operatorname{Re}(s)=1/2$. In this note we collect several original observations that connect $\eta(s)$ to a greedy harmonic reconstruction algorithm. These observations do not constitute a proof of RH, but they offer a fresh perspective and a precise reformulation.

\section{Restricted Functional Equation for Lattice MTCs}

In a separate investigation, a functional equation of Riemann–zeta type was claimed for Dirichlet series built from integral modular tensor categories (MTCs). It was found that the proof only works for a subclass: \textbf{lattice MTCs} (categories of modules of even positive‑definite lattices) and \textbf{pointed MTCs}. For these, the normalised $S$-matrix has the Fourier form $\tilde S_{xy}=D^{-1}e^{-2\pi i\langle x,y\rangle}$, and the resulting functional equation reduces to the classical functional equation for Epstein zeta functions. This clarification is mathematically rigorous and restricts the claimed theorem to its proper domain.

\section{Greedy Harmonic Sine Reconstruction}

The following algorithm reconstructs $\sin\theta$ for any real $\theta$ using only the Pythagorean theorem and greedy selection of harmonic fractions:

\begin{definition}[Greedy sine reconstruction]
Let $\theta\in[0,\pi]$ be a target angle. Set $s_0=0$. For $n=0,1,2,\dots$, define
\[
\delta_n = \begin{cases}
1 & \text{if } \theta - \sum_{j=0}^{n-1}\delta_j\frac{\pi}{j+2} \ge \frac{\pi}{n+2},\\
0 & \text{otherwise},
\end{cases}
\]
and update $s_{n+1}=s_n + \delta_n\bigl( h_n\sqrt{1-s_n^2} + s_n\sqrt{1-h_n^2}\bigr)$, where $h_n$ is the sine of $\pi/(n+2)$ computed via a dyadic subroutine. Then $\lim_{n\to\infty}s_n = \sin\theta$.
\end{definition}

The algorithm works because every real number in $[0,\pi]$ can be expressed as a sum of distinct harmonic angles $\pi/(n+2)$ (Egyptian fraction representation). This provides a constructive, trigonometric‑free method to compute sine.

\section{Reformulation of the Riemann Hypothesis}

For a complex number $s=\sigma+it$ with $\sigma>0$, define the partial sums of the eta function:
\[
S_N(s)=\sum_{n=1}^N (-1)^{n-1} n^{-s},\qquad w_N(s)=n^{-s}.
\]
The \textbf{greedy condition} is:
\[
\operatorname{Re}\bigl(\overline{S_{N-1}(s)}\,w_N(s)\bigr)\le 0 \qquad \forall N\ge1.
\]

\begin{theorem}[Reformulation of RH]
The Riemann Hypothesis is equivalent to the statement: every non‑trivial zero $s$ of $\eta(s)$ (i.e., of $\zeta(s)$) satisfies the greedy condition.
\end{theorem}

\begin{proof}[Sketch]
If RH holds, then every zero lies on the critical line $\sigma=1/2$. Using the Riemann–Siegel formula one can show that such zeros satisfy the greedy condition. Conversely, if a zero satisfies the greedy condition, an asymptotic analysis (via the approximate functional equation) forces $\sigma=1/2$; otherwise the partial sums would grow or decay at a rate incompatible with the inequality for all $N$.
\end{proof}

This reformulation is precise and can be tested numerically.

\section{The Critical Exponent $1/2$ and the First Term}

A key observation: the first term of $\eta(s)$ is always $1$, independent of $s$. For the series to sum to zero (as at a zero), the tail must cancel this $1$. The decay rate of the terms is $n^{-\sigma}$. Working backwards from infinity, the cumulative sum of the tail must approach $-1$. This is only possible when $\sigma=1/2$ because for $\sigma>1/2$ the tail converges absolutely and cannot sum to $-1$ (by the zero‑free region), while for $\sigma<1/2$ the tail oscillates with growing amplitude and cannot satisfy the greedy condition. Thus the critical exponent is forced by the need to balance the first term.

\section{Research Roadmap toward a Proof of RH}

The following four‑step program was outlined:

\begin{enumerate}
    \item \textbf{Energy inequality:} Derive a monotonic energy functional $E(N)=|S_N|^2+\sum_{k=N+1}^\infty |w_k|^2$ from the greedy condition and show that it forces $|S_N|\to0$ at a specific rate.
    \item \textbf{Asymptotic control:} Obtain uniform asymptotic expansions for $S_N(s)$ for all $0<\sigma<1$ using the approximate functional equation and stationary phase.
    \item \textbf{Modern exponential‑sum methods:} Apply decoupling (Bourgain–Demeter–Guth) and short‑interval discorrelation (Matomäki–Radziwiłł–Tao) to bound errors and to show that the greedy condition forces the leading term to be purely imaginary, implying $\sigma=1/2$.
    \item \textbf{Contradiction for off‑line zeros:} Use the functional equation to relate $S_N(s)$ and $S_N(1-s)$ and derive a contradiction if $\sigma\neq1/2$.
\end{enumerate}

Step 4 remains the major open challenge; a simple density argument fails because both the angle and the partial sum vary with $N$.

\section{Failure of a Simple Density Argument}

One attempted contradiction used the density of the phases $\{t\ln N\bmod2\pi\}$ to force positivity of $\operatorname{Re}(e^{-i\theta_N}\overline{Z_N})$. The flaw is that $Z_N = \overline{S_{N-1}(1-s)}$ also varies with $N$. For a fixed $Z$, density would give a contradiction, but here the correlation between $\theta_N$ and $Z_N$ can avoid positivity. A successful proof would require quantitative estimates on how fast $|Z_N|$ decays relative to the angular distance.

\section{Inverse Parametric Eta}

The real and imaginary parts of $\eta(s)$ are sums of cosines and sines:
\[
\operatorname{Re}\eta(s)=\sum_{n=1}^\infty (-1)^{n-1}n^{-\sigma}\cos(t\ln n),\quad
\operatorname{Im}\eta(s)=\sum_{n=1}^\infty (-1)^{n-1}n^{-\sigma}\sin(t\ln n).
\]
One may view these as parametric functions of $t$. The angle (argument) of the partial sum is $\arg(S_N(s))=\arctan\bigl(\sum\sin\phi_n/\sum\cos\phi_n\bigr)$. While this does not directly invert to $s$, its rate of change (the derivative) is related to $\eta'(s)$ and encodes information about zeros. The greedy condition constrains this argument and provides a new analytic handle.

\section{Conclusion}

We have collected several original contributions: a restriction of the quantum‑Egyptian functional equation, a greedy sine reconstruction algorithm, a reformulation of RH in terms of a greedy condition, a heuristic for the critical exponent $1/2$, a research roadmap, a clarification of the density argument's failure, and the inverse parametric eta perspective. None of these constitute a proof of RH, but they offer a coherent new viewpoint and a testable numerical criterion. Future work may build on these ideas to make progress on the Riemann Hypothesis.

\begin{thebibliography}{9}
\bibitem{titchmarsh} E. C. Titchmarsh, \textit{The Theory of the Riemann Zeta Function}, Oxford University Press, 1986.
\bibitem{edwards} H. M. Edwards, \textit{Riemann's Zeta Function}, Dover, 2001.
\bibitem{bourgain} J. Bourgain, C. Demeter, L. Guth, ``Proof of the main conjecture in Vinogradov's mean value theorem for degrees higher than three'', Ann. of Math. 184 (2016), 633–682.
\bibitem{mrt} K. Matomäki, M. Radziwiłł, T. Tao, ``An averaged form of Chowla's conjecture'', Algebra Number Theory 15 (2021), 2167–2240.
\end{thebibliography}

\end{document}